\documentclass{article} %210 mm × 297 mm
\usepackage[utf8]{inputenc}
\usepackage[a4paper,left=1.5cm, right=1.5cm, top=2cm, bottom=2cm]{geometry}
\usepackage{graphicx}
%\usepackage{blindtext}
\usepackage{multirow}
\usepackage{mhchem}
\usepackage{url}
\usepackage{subfigure}
\usepackage{amsfonts,latexsym} % para tener disponibilidad de diversos simbolos
\usepackage{enumerate}
%\usepackage{booktabs}
\usepackage{float}
%\usepackage{threeparttable}
\usepackage{array,colortbl}
%\usepackage{ifpdf}
%\usepackage{rotating}
\usepackage{cite}
\usepackage{stfloats}
\usepackage{url}
\usepackage{listings}
\usepackage{fancyhdr}

\usepackage{color}
%\usepackage{hyperref}
%\usepackage{wrapfig}
\usepackage{array}
%\usepackage{multirow}
%\usepackage{adjustbox}
%\usepackage{nccmath}
\usepackage[dvipsnames]{xcolor}


\newcommand{\titulo}{Abgabe 1; Diskrete Mathematik}
\newcommand{\nombre}{Liliana Sanfilippo}
\newcommand{\FCE}{\textbf{WiSe 2024/2025}}

 

\usepackage{fancyhdr}
\pagestyle{fancy}
\fancyhead{}
\fancyfoot{}
\fancyhead[LO]{\small\nombre}
\fancyhead[RE]{\small\titulo}
\fancyhead[LO]{\small\FCE}
\fancyfoot[RO,LE] {\thepage}

\setcounter{page}{1} %Número de la página, la editorial lo cambiará



\usepackage[onehalfspacing]{setspace}
\begin{document}


\begin{tabular}{p{12cm}}
{{\Large\bf\titulo}}\\
\vspace{0.2cm}\\
 {\large\nombre}\\
\end{tabular}
\vspace{0.2cm}\\
\section{Präsenzaufgaben}



\subsection{Aufgabe 1 - \textcolor{blue}{Teilmengen}}
Im Hof spielen 20 Kinder. Ein Elternteil hat 5 Schokoriegel zu verschenken. Wie viele Möglichkeiten gibt es, wenn
\begin{itemize}
    \item alle Schokoriegel verschieden sind, jedes Kind darf maximal einen haben? \\ 
    \textcolor{blue}{\textbf{Definition 2.1.} Die Anzahl der $r$-elementigen Teilmengen einer $n$-elementigen Menge wird mit $\binom{n}{r}$ bezeichnet.
} \newline
    \textcolor{blue}{jedes Kind darf maximal einen haben $\rightarrow$ Anzahl der r-elementigen Teilmengen einer n-elementigen Menge, wobei n Anzahl der Kinder und r Anzahl der Schokoriegel: $\binom{n}{r} = \binom{20}{5}$ } \\
\textcolor{blue}{alle Schokoriegel verschieden $\rightarrow$ Reihenfolge der Verteilung der Schokoriegel ist relevant: $5*4*3*2*1 = 120 = 5!$ Möglichkeiten } \\
    \textcolor{blue}{Daher: $\binom{20}{5}  * 5!$ = $\frac{20!}{15!}$ }
 \item alle Schokoriegel gleich sind, jedes Kind darf maximal einen haben? \\
 \textcolor{blue}{\textbf{Definition 2.1.} Die Anzahl der $r$-elementigen Teilmengen einer $n$-elementigen Menge wird mit $\binom{n}{r}$ bezeichnet.
} \\
 \textcolor{blue}{jedes Kind darf maximal einen haben $\rightarrow$ Anzahl der r-elementigen Teilmengen einer n-elementigen Menge, wobei n Anzahl der Kinder und r Anzahl der Schokoriegel: $\binom{n}{r} = \binom{20}{5}$ } \\
\textcolor{blue}{alle Schokoriegel gleich $\rightarrow$ Reihenfolge der Verteilung der Schokoriegel ist irrelevant } \\
    \textcolor{blue}{Daher: $\binom{20}{5}$ }
 \item alle Schokoriegel gleich sind, die Kinder nehmen wie viel sie wollen/können? \\
 \textcolor{blue}{\textbf{Satz 2.9.} Es gibt $\binom{n+r-1}{r}$ Möglichkeiten, $r$ Elemente aus einer $n$-elementigen Menge mit Wiederholung auszuwählen.
}\\ 
\textcolor{blue}{$\binom{n+r-1}{r}$ = $\binom{20+5-1}{5}$ = $\binom{24}{5}$}
 \item alle Schokoriegel verschieden sind, die Kinder nehmen wie viel sie wollen/können? \\
 \textcolor{blue}{\textbf{Potenzregel 1.5.} \quad $|X^n| = |X| \cdot \cdots \cdot |X| = |X|^n$
} \\
\textcolor{blue}{Also: $20^5$}
\end{itemize}


\subsection{Aufgabe 2}
In einem Modul muss jede*r der 35 Studierenden eines Jahrgangs genau drei von fünf
möglichen Veranstaltungen wählen. In der ersten Veranstaltung nehmen 18 teil, in der
zweiten 23, in der dritten 20 und in der vierten 27.
\\ \newline
Wie viele belegen die fünfte Veranstaltung?
\\ 
\color{blue}
\begin{itemize}
    \item Studenten in Veranstaltung 1: 18 
    \item Studenten in Veranstaltung 2: 23
    \item Studenten in Veranstaltung 3: 20
    \item Studenten in Veranstaltung 4: 27 
    \item Studenten in Veranstaltung 5: x
\end{itemize}
Jeder muss drei Veranstaltungen belegen, also sollten 105 Einschreibungen in ein Modul vorliegen. 
\begin{equation*}
    18 + 23 + 20 + 27 + x = 105
\end{equation*}
\begin{equation*}
    88 + x = 105
\end{equation*}
\begin{equation*}
     x = 17
\end{equation*}
In Veranstaltung 5 befinden sich 17 Studenten.
\color{black}

\subsection{Aufgabe 3}
Sicherheitsschlüssel haben Kerben an genau bestimmten Positionen. Wenn die Kerben
acht verschiedene Tiefen haben können, wie viele Positionen muss man dann vorsehen,
damit man eine Billion verschiedene Schlüssel herstellen kann?
\color{blue}
Anzahl Tiefen: 8 \\
Anzahl Kerben: x \\
\begin{equation*}
    8^x = 1.000.000.000.000
\end{equation*}
\begin{equation*}
    log(8^x) = log(1.000.000.000.000)
\end{equation*}
\begin{equation*}
   x * log(8) = log(1.000.000.000.000)
\end{equation*}
\begin{equation*}
   x * log(8) = log(10^{12})
\end{equation*}
\begin{equation*}
   x * log(8) = 12
\end{equation*}
\begin{equation*}
   x  = \frac{12}{log(8)} \approx 13.29
\end{equation*}
Aufrunden, damit man tatsächlich eine Billion hinbekommt: 14
\color{black}


\subsection{Aufgabe 4}
Es sei $X = \{1, 2, \dots, 100\}$ und $A \subseteq X$.
\begin{itemize}
    \item[(a)] Zeigen Sie: Wenn $|A| = 55$, dann gibt es $a, b \in A$ mit $a - b = 9$.
    \color{blue}

\color{black}
    \item[(b)] Gilt dies auch für $|A| = 54$?
\end{itemize}
\color{blue}

\color{black}


\section{Aufgaben zur Abgabe}
\subsection{Aufgabe 1 (4 Punkte)}
Im Parlament eines Landes gibt es 151 Sitze und drei Parteien. Wie viele Möglichkeiten
der Sitzverteilung gibt es, so dass keine Partei eine absolute Mehrheit hat?
\color{blue}

\color{black}

\subsection{Aufgabe 2 (4 Punkte)}
Die acht Professor*innen A, B, C, . . . , H einer Fakultät haben mehrere Kommissionen
mit je vier Mitgliedern gebildet. Jede*r Professor*in gehört genau drei Kommissionen an
und keine zwei Kommissionen sind gleich zusammengesetzt. Wie viele Kommissionen gibt
es? Begründen Sie Ihre Antwort und finden Sie eine mögliche Zusammenstellung der
Kommissionen.
\color{blue}

\color{black}

\subsection{Aufgabe 3 (4 Punkte)}
Zeigen Sie: Unter je fünf Punkten der Ebene mit ganzzahligen Koordinaten gibt es zwei,
deren Mittelpunkt ebenfalls ganzzahlige Koordinaten hat.
\color{blue}

\color{black}

\subsection{Aufgabe 4 (4 Punkte)}
Auf wie viele Arten kann ein König von der linken unteren Ecke eines Schachbrettes nach
der rechten oberen ziehen, wenn er stets nach oben, nach rechts oder diagonal nach rechts
oben zieht?
\color{blue}

\color{black}


\end{document}